\chapter{Discussion et Perspectives}
\label{ch:discussion_perspectives}

\section*{Introduction}
\phantomsection
\addcontentsline{toc}{section}{Introduction}

Ce chapitre revient sur ce qui a été
accompli concrètement pendant ce projet,
les limites réelles observées pendant
les tests, et ce qui resterait à faire
avant un déploiement urbain réel.

\section{Synthèse des contributions}

Ce travail a produit quatre contributions
complémentaires, dont la valeur ne réside
pas dans chaque brique prise isolément
mais dans leur intégration dans un seul
prototype mesuré sur Raspberry Pi 5.

\subsection{Pipeline embarqué complet}

Le système déploie six modules
interconnectés sur un seul Raspberry
Pi 5, sans dépendance cloud : acquisition
vidéo, détection YOLO26n, anonymisation
Privacy-by-Design, tracking IoU par zones
polygonales, décision MDP et dashboard
Flask. Le fait que ces six modules
fonctionnent ensemble sur un CPU ARM à
ressources limitées, avec les compromis
de vitesse et de précision documentés,
constitue la contribution principale.

\subsection{Dataset hybride méthodique}

Le dataset de 6 761 images a été
construit selon une méthodologie
rigoureuse en trois étapes — split
sans fuite de données, renforcement
CityPersons, crops tiny persons —
pour répondre à une contrainte que
les datasets publics disponibles ne
couvraient pas : des scènes de
surveillance d'intersection en angle
fixe, avec des classes adaptées au
contexte marocain. La méthodologie
est documentée et reproductible.

\subsection{Benchmark empirique sur Pi 5}

Le benchmark comparatif de six
architectures YOLO en conditions
réelles sur Raspberry Pi 5 produit
des données empiriques absentes de
la littérature pour cette combinaison
matérielle. Deux résultats sont
directement utiles à d'autres travaux :
les variantes small sont incompatibles
avec le temps réel sur CPU ARM (moins
de 3 FPS dans tous les cas), et la
quantification INT8 dynamique ralentit
l'inférence sur ONNX Runtime CPU au
lieu de l'accélérer.

\subsection{Privacy-by-Design conforme
            Loi 09-08}

Positionner le module d'anonymisation
avant le tracking, avec un tracker IoU
qui n'utilise que les coordonnées des
boîtes, permet de respecter la Loi 09-08
sans dégrader les performances
fonctionnelles du pipeline. Le rejet
motivé de ByteTrack pour raisons
juridiques et matérielles offre un
modèle d'architecture reproductible
pour d'autres systèmes de surveillance
urbaine en contexte marocain.

\section{Bilan des hypothèses}

Les résultats du Chapitre~\ref{ch:resultats}
sont résumés ici avec les nuances
importantes :

\begin{table}[H]
\centering
\caption{Bilan de validation des
         hypothèses de recherche}
\small
\begin{tabularx}{\textwidth}{cXXX}
\toprule
\textbf{H} &
\textbf{Hypothèse} &
\textbf{Résultat clé} &
\textbf{Statut} \\
\midrule
H1 &
Détection en conditions urbaines &
0.792 mAP50 test &
Validée — limite person
documentée \\
H2 &
Dataset hybride &
6 761 images, gain person
800→960, sans ablation
par source &
Validée empiriquement \\
H3 &
Edge embarqué &
Nano : 7.6 FPS ;
Step 3 800 stride 2 :
7.84 frames source/s &
Validée nano ;
partielle Step 3 \\
H4 &
Anonymisation &
100\% détections person
floutées avant tracking &
Validée — risque résiduel
sur non-détections \\
H5 &
Décision MDP &
3 scénarios NS/EW/Emergency
produits correctement &
Validée sur phases
implémentées \\
\bottomrule
\end{tabularx}
\label{tab:hypotheses_bilan_ch6}
\end{table}

\section{Comparaison finale avec
         la littérature}

Le tableau suivant met à jour la synthèse
comparative du Chapitre~\ref{ch:etat_art}
avec les résultats réels mesurés dans
ce travail.

\begin{table}[H]
\centering
\caption{Comparaison finale —
         résultats mesurés vs littérature}
\scriptsize
\begin{tabularx}{\textwidth}{lXXXXX}
\toprule
\textbf{Travail} &
\textbf{Modèle} &
\textbf{Hardware} &
\textbf{mAP50} &
\textbf{FPS réel} &
\textbf{Privacy} \\
\midrule
Hazarika (2024) &
tiny-YOLOv3 &
Raspberry Pi &
— &
$\sim$5 &
Non \\
Sapkota (2025) &
YOLO26n &
Jetson Orin &
— &
$>$30 &
Non \\
Krishna (2026) &
YOLOv8 &
Edge device &
— &
— &
Non \\
Mahalakshmi (2025) &
Deep Learning &
PC standard &
— &
— &
Non \\
\midrule
\rowcolor{PrimaryBlue!10}
\textbf{Ce travail} &
\textbf{YOLO26n} &
\textbf{Pi 5 8GB} &
\textbf{79.2\%} &
\textbf{7.6 (nano)} &
\textbf{PbD} \\
\rowcolor{PrimaryBlue!5}
\textit{Step 3 960} &
\textit{YOLO26n} &
\textit{Pi 5 8GB} &
\textit{79.2\%} &
\textit{2.16 (analyse)} &
\textit{PbD} \\
\bottomrule
\end{tabularx}
\vspace{0.3em}
\begin{flushleft}
\scriptsize\textit{— : métrique non
rapportée pour ce hardware.
PbD : Privacy-by-Design conforme
Loi 09-08. FPS mesuré sur le
dispositif cible de chaque travail.
Le nano valide le temps réel ;
Step 3 960 est le mode haute précision.}
\end{flushleft}
\label{tab:comparaison_finale}
\end{table}

Ce tableau illustre le positionnement
de ce travail. Hazarika (2024) est le
seul autre travail déployé sur Raspberry
Pi, mais avec une architecture plus
ancienne et sans Privacy-by-Design.
Sapkota (2025) utilise YOLO26n mais
sur Jetson Orin — hardware dix fois
plus puissant. Notre contribution
documente pour la première fois, à
notre connaissance, les performances
mesurées de YOLO26n sur Raspberry Pi 5
avec anonymisation intégrée et décision
adaptative, dans un contexte de smart
city marocaine.

\section{Limites du système}

\subsection{Dataset insuffisamment
            représentatif}

Les 133 images collectées localement
à Lqliaa restent insuffisantes pour
couvrir la diversité des conditions
de trafic marocain. Les conditions
nocturnes, la pluie, les intersections
à géométrie complexe et les véhicules
atypiques — motos chargées, véhicules
à traction animale encore présents
dans certains quartiers — ne sont pas
représentés. Un modèle entraîné sur
ce dataset peut ne pas généraliser
correctement à ces situations.

\subsection{Détection des petits piétons}

La classe person reste la plus difficile
du système. Un rappel de 0.353 sur le
split test signifie qu'environ un tiers
des piétons réellement présents dans
les frames de test ne sont pas détectés
— et donc pas anonymisés. C'est le
principal risque juridique résiduel
vis-à-vis de la Loi 09-08. Les efforts
sur CityPersons, les tiny crops, la
résolution 960 et les ROI ont amélioré
ce chiffre, mais pas suffisamment pour
garantir une couverture complète.

\subsection{Décision MDP sans dynamique
            temporelle}

La politique de scoring pondéré réagit
à l'état instantané de l'intersection.
Elle ne tient pas compte de ce qui
s'est passé dans les minutes précédentes,
ni de ce qui risque de se passer dans
les minutes suivantes. Un axe avec une
densité croissante ne bénéficie d'aucune
anticipation — la décision attend que
la congestion soit déjà présente pour
réagir. C'est une limite connue des
approches déterministes par rapport aux
approches par apprentissage par
renforcement.

\subsection{Compromis précision/latence
            sur Pi 5}

Le modèle Step 3 960 atteint la meilleure
précision disponible pour ce contexte,
mais son débit de 2.16 FPS avec ROI sur
Pi 5 le réserve à un mode d'analyse.
Même sans ROI, il n'atteint que 2.51 FPS.
L'export 800 améliore le débit à 3.94 FPS,
et le mode stride 2 couvre 7.84 frames
source par seconde — mais en n'inférant
qu'une frame sur deux. Pour un déploiement
multi-intersections à l'échelle d'une
ville, ce compromis ne serait pas
suffisant sans accélération matérielle
ou réduction de la fréquence d'inférence.

\section{Perspectives}

\subsection{Améliorer la couverture
            des piétons}

C'est la priorité la plus directement
liée à la conformité légale. Deux pistes
sont envisageables. La première est
un modèle dédié aux petits piétons,
plus léger qu'un modèle de détection
généraliste, activé uniquement sur
les zones piétonnes — ce qui limiterait
son impact sur le débit. La seconde
est l'exploration de techniques
d'anonymisation irréversibles comme
ReCAM~\cite{jeremiah2024}, qui
permettraient de traiter les personnes
partiellement détectées avec plus
de robustesse.

\subsection{Implémenter les phases
            Pedestrian et All-Red}

La politique MDP du Chapitre~\ref{ch:conception}
formalise cinq phases, dont Pedestrian
et All-Red. Le système actuel n'implémente
pas ces deux règles dans le code de
décision. Les ajouter permettrait de
compléter la validation H5 et d'améliorer
la gestion des intersections à forte
fréquentation piétonne.

\subsection{Vers un agent DRL léger}

Remplacer la politique déterministe par
un agent de Deep Reinforcement Learning
distillé — entraîné en simulation puis
compressé pour tenir sur Pi 5 — pourrait
améliorer les décisions sur des scénarios
de trafic complexes et dynamiques.
C'est une perspective à plus long terme,
qui nécessite un simulateur routier
représentatif du contexte marocain.

\subsection{Coordination multi-intersections}

Le système actuel traite une seule
intersection de façon indépendante.
Une coordination entre plusieurs
intersections via MQTT permettrait
d'optimiser les green waves sur des
axes entiers — comme l'Avenue Mohammed VI
à Agadir. C'est une extension naturelle
de l'architecture actuelle, qui ne
nécessite pas de refondre le pipeline.

\subsection{Dataset marocain public}

Publier le dataset Lqliaa enrichi sous
une licence ouverte comblerait un manque
réel dans la littérature sur le trafic
urbain marocain. C'est une contribution
durable qui dépasse le cadre de ce
mémoire.

\subsection{Déploiement pilote à Agadir}

Le programme Agadir Smart City offre
un cadre institutionnel pour un
déploiement pilote sur une intersection
réelle. Un test de 30 jours sur
l'intersection de Lqliaa, avec mesure
des temps d'attente, du débit et de la
stabilité des phases, permettrait de
valider les performances en conditions
réelles — ce que les vidéos Bellevue ne
peuvent pas faire.

\section*{Synthèse du chapitre}
\addcontentsline{toc}{section}{Synthèse}

Ce chapitre a présenté quatre contributions
concrètes, un bilan honnête des hypothèses
validées avec leurs limites, et six
perspectives organisées par ordre de
priorité. Les deux points les plus
urgents avant un déploiement réel sont
l'amélioration de la couverture des
piétons et l'implémentation des phases
Pedestrian et All-Red dans le module M5.